\chapter{Conclusions \& Outlook}

\section{Summary of results}

This thesis investigated the use of nuclear magnetic resonance for searches for axionlike dark-matter particles through the nuclear-gradient coupling, focusing on the CASPEr-Gradient-Low-Field experiment at Mainz.

The theoretical analysis in Chapter~\ref{chap:theory} and in Sections~\ref{sec:experiment:single}--\ref{sec:experiment:optimization} of Chapter~\ref{chap:experiment} (adapted from Ref.\,\cite{Yuzhe2023_FrequencyScanning}) established the following key results:

\begin{itemize}
    \item \textbf{Optimal step size.}
    The frequency-step size for scanning should equal the sensitive bandwidth of a single measurement, namely the larger of the NMR linewidth $\Gamma_n$ and the ALP linewidth $\Gamma_a$.

    \item \textbf{Optimal relaxation times.}
    The transverse relaxation times $T_2$ and $T_2^*$ should be maximized.
    Longer $T_2^*$ is beneficial in all regimes: it increases either the r.m.s.\ tipping angle or the fraction of on-resonance spins.
    Longer $T_2$ improves sensitivity when $T_2 \ll \tau_a$.

    \item \textbf{Sensitivity scaling.}
    The sensitivity to $|g_{\mathrm{aNN}}|$ scales as $T_{\mathrm{tot}}^{-1/4}$ with total measurement time in all regimes considered ($T \gg \tau_a, T_2^*, T_2$).
    This slow scaling means improving $T_2$, $T_2^*$, or the detector noise $\sigma_0$ is more efficient than simply accumulating more measurement time.

    \item \textbf{Optimal dwell-time allocation.}
    For regimes~(1)--(3), the dwell time should be distributed uniformly across frequency steps ($T \propto \nu^0$).
    For regime~(4) ($T_2^*=T_2 \ll \tau_a$), the optimal strategy is $T \propto \nu^{-1}$.

    \item \textbf{Four regimes.}
    The parameter space is divided into four regimes by the relative magnitudes of $T_2$, $T_2^*$, and $\tau_a$.
    Each regime has a distinct signal amplitude and optimal dwell time, as summarized in Table~\ref{tab:summary_of_regimes}.
\end{itemize}

These results were applied in a proof-of-principle ALP search with the CASPEr-Gradient-LF apparatus (adapted from Ref.\,\cite{Walter2024_LFthermal}), described in Sections~\ref{sec:results:measurement}--\ref{sec:results:exclusion} of Chapter~\ref{chap:results}:

\begin{itemize}
    \item \textbf{First laboratory search in the target mass range.}
    A thermally polarized methanol sample was used to search for the ALP-proton gradient coupling in the mass range $5.576741$--$5.577733\,\si{\neV}/c^2$ (Compton frequency \SI{240}{\Hz} around \SI{1.348570}{\MHz}).
    This is the first dedicated laboratory search in this mass range.

    \item \textbf{No dark-matter candidates found.}
    Three sequential scans were performed.
    No ALP signal candidates survived the cross-scan consistency check, yielding the constraint $|g_\mathrm{ap}| \lesssim \SI{3e-2}{\GeV^{-1}}$ at \SI{95}{\%} confidence level.

    \item \textbf{Methodology validated.}
    The data-analysis pipeline---Savitzky-Golay spike suppression, baseline removal, matched filtering, and NPE hypothesis test---was validated through synthetic-signal injection.
    The overall analysis efficiency is $\eta_\mathrm{rec} = \SI{88.2\pm4.6}{\%}$.

    \item \textbf{Sensitivity limited by thermal polarization.}
    The proton thermal polarization at \SI{180}{\kelvin} and \SI{31.67}{\milli\tesla} is only $\approx\num{1.8e-7}$, which dominates the sensitivity budget.
    The apparatus is otherwise ready for more sensitive searches with hyperpolarized samples.
\end{itemize}


\section{Outlook}

The framework developed here can be extended in several directions:

\begin{itemize}
    \item \textbf{Hyperpolarized samples.}
    Hyperpolarized \ch{^{129}Xe} or methanol can increase the polarization by six or more orders of magnitude relative to thermal polarization, directly improving $M_0$ and hence the sensitivity prefactor $\eta$ in Equation~\eqref{eq:eta}.
    The $T_1$ relaxation of hyperpolarized samples limits the measurement window and modifies the optimal dwell-time strategy.
    Examples of projected sensitivities are given in Table~\ref{tab:parameters}.

    \item \textbf{Seeded measurements.}
    Applying a transverse oscillating field at the resonance frequency (``seeding'') to heterodyne the ALP signal was outside the scope of this work.
    Seeding could improve sensitivity by coherently amplifying the ALP-induced tipping angle \cite{Aybas_2021_Quantum_sensitivity}.

    \item \textbf{Improved field homogeneity and longer $T_2$.}
    Reaching the design inhomogeneity of \SI{2}{\ppm} over the sample volume would lengthen $T_2^*$ and increase the fraction of on-resonance spins.
    Combined with samples having long intrinsic $T_2$, this improvement is beneficial in all four regimes.

    \item \textbf{High-field extension.}
    The CASPEr-Gradient-High-Field experiment will operate with a \SI{14.1}{\tesla} magnet, reaching Larmor frequencies up to \SI{600}{\MHz} and accessing higher ALP masses.
    The scanning framework derived here applies directly to that regime.

    \item \textbf{Low-frequency regime.}
    Below about \SI{1}{\kHz}, $\tau_a$ can exceed \SI{1000}{\second}, and the assumption $T \gg \tau_a$ breaks down.
    An adaptive scanning strategy that accounts for the long coherence time will be required for CASPEr-ZULF \cite{CASPEr_ZULF_2019} and related low-frequency experiments.
\end{itemize}
